\newpage
\section{NNPDF3.1 全局分析}

\label{sec:results}

本节给出 NNPDF3.1 在 LO、NLO 与 NNLO 的全局分析结果，并与前一版本
NNPDF3.0 以及其他近期 PDF 集合比较。
%
这里给出的结果采用第~\ref{sec:expdata} 节讨论的表~\ref{tab:completedataset}--\ref{tab:completedataset3}
完整数据集。对单个测量影响的讨论将与基于缩减数据集的 PDF 确定一起在第~\ref{sec:impactnewdata} 节进行。

在简要的方法学小结之后，
我们讨论拟合质量，
继而逐一考察各 PDF 及其不确定度。
我们将 NNPDF3.1 PDF 与 NNPDF3.0 以及 CT14~\cite{Dulat:2015mca}、MMHT2014~\cite{Harland-Lang:2014zoa}
和 ABMP16~\cite{Alekhin:2017kpj} 比较。
接着我们考察独立参数化粲的影响——这是 NNPDF3.1 最主要的方法学改进。最后我们讨论理论不确定度，既包括与 QCD 参数相关的，也包括 PDF 确定所用理论中缺失的高阶修正。

本节中所有 NLO 与 NNLO 的 NNPDF3.1 结果都用 CMC~\cite{Carrazza:2015hva}
优化的 100 副本蒙特卡罗集合产生（见下文第~\ref{sec:delivery} 节）：尽管只含 100 个副本，这些集合重现了至少约 400 副本集合的统计特征（见第~\ref{sec:compression} 节）。此处仅展示部分结果：更全面的结果可从公共存储库获得，见第~\ref{sec:delivery} 节。

\subsection{方法学}
\label{sec:methodology}

NNPDF3.1 PDF 的确定方法与 NNPDF3.0 基本相同：
唯一的重要变化是如今粲被独立参数化。PDF 参数化与文献~\cite{Ball:2014uwa}
第 3.2 节所述相同，包括预处理（preprocessing）的处理方式，只是该文献式 (3.4) 的 PDF
基底现在补充了一个额外的粲 PDF，它像其他所有 PDF 一样参数化（依文献~\cite{Ball:2016neh} 式 (2)）。
凡粲 PDF 被独立参数化时，PDF 都在标度 $Q_0=1.65$~GeV 处参数化。为便于比较，我们也提供微扰生成粲构造的 PDF 集合；
这些集合中 PDF 在 $Q_0=1.0$~GeV 处参数化。这保证当粲被独立参数化时参数化标度总高于粲质量，
而当粲为微扰生成时则低于粲质量。

与文献~\cite{Ball:2014uwa}一样，全文默认取 $\alpha_s(m_Z)=0.118$，但也对若干不同 $\alpha_s$ 值进行了拟合（见第~\ref{sec:delivery} 节）。
全文采用重夸克极点质量，主要理由是：对用于 PDF 确定的单举观测量而言 $\overline{\rm MS}$ 质量并不合适，因为它会在阈值区扭曲微扰展开~\cite{Ball:2016qeg}。
重夸克极点质量的默认值按希格斯截面工作组的建议取为：粲 $m_c=1.51$~GeV、底 $m_b=4.92$~GeV~\cite{deFlorian:2016spz}；
对应于文献~\cite{deFlorian:2016spz} 给出的 $\pm 1$ 倍标准差不确定度带的不同粲质量值对应的 PDF 集合也已提供，见第~\ref{sec:delivery} 节。


\subsection{拟合质量}
\label{sec:fquality}

%%%%%%%%%%%%%%%%%%%%%
\begin{table}[h]
\centering
\includegraphics[page=1,width=0.98\textwidth]{images/C.pdf}
\caption{\small 全局拟合以及 NNPDF3.1 LO、NLO 与 NNLO PDF 拟合中全部数据集的 $\chi^2/N_{\rm dat}$ 数值；同时给出用 NNPDF3.0 NLO 与 NNLO PDF 得到的数值：方括号中的数字对应 NNPDF3.0 未参与拟合的数据。注意 NNPDF3.0 数值采用 NNPDF3.1 理论设置计算，因此比文献~\cite{Ball:2014uwa} 所引数值略差。
\label{tab:chi2tab_31-nlo-nnlo-30}}
\end{table}


%%%%%%%%%%%%%%%%%%%%%


表~\ref{tab:chi2tab_31-nlo-nnlo-30} 给出全局拟合以及 NNPDF3.1 LO、NLO 与 NNLO PDF 拟合中各数据集各自的 $\chi^2/N_{\rm dat}$ 数值，并与 NNPDF3.0 NLO、NNLO 的相应值比较。$\chi^2$
用包含全部关联的协方差矩阵计算，即各实验发表的协方差矩阵。
考察该表可见拟合质量从 LO 到 NLO 到 NNLO 逐步改善：不仅 LO 与 NLO 之间有显著改进，从 NLO 到 NNLO 也有明显改善。值得注意这在 NNPDF3.0 中并非如此——那里 NNLO 的拟合质量实际上比 NLO 略差（见文献~\cite{Ball:2014uwa} 表 9）。这反映了 NNPDF3.1 所含强子过程比例的提高（其 NNLO 修正往往相当可观），也可能部分源于方法学改进。

NNPDF3.1 的总体拟合质量明显好于用 NNPDF3.0 PDF 得到的结果。
这对未纳入 NNPDF3.0 的 LHC 测量当然在预料之中；但有意思的是，3.0 中已有的 HERA 测量（尽管当时是略有不同的未合并形式）现在也拟合得更好。
不过，旧版 NNPDF3.0 PDF 对全部新数据的描述质量仍相当可以接受，
表明 NNPDF3.0 与 NNPDF3.1 总体上是相容的。
注意表~\ref{tab:chi2tab_31-nlo-nnlo-30} 中 NNPDF3.0 数值是用 NNPDF3.1 理论设置计算的，特别是重夸克质量不同于 NNPDF3.0 PDF 拟合所用数值。因此文献~\cite{Ball:2014uwa} 表 9 所示的 NNPDF3.0 拟合质量略好于表~\ref{tab:chi2tab_31-nlo-nnlo-30} 所示者，但即便如此 NNPDF3.1 的拟合质量仍更佳。具体到 HERA 数据，采用一致理论设置的 NNPDF3.0 拟合质量可从文献~\cite{Czakon:2016olj} 表 7 读出：其值为 $\chi^2/N_{\rm dat}=1.21$，从而确实表明 NNPDF3.1 给出了更好的描述。这一改进的原因将在下文第~\ref{sec:results-mc} 节讨论。

对许多新的 LHC 测量，只有在 NNLO 才能实现对数据的良好描述。
ATLAS、CMS 与 LHCb 实验的总 $\chi^2/N_{\rm dat}$ 在 NNLO 分别为 1.09、1.06 与 1.47，
而 NLO 时分别为 1.36、1.20 与 1.62。从 NLO 到 NNLO 改善最大的数据集正是实验不确定度最小的那些。例如
ATLAS $W,Z$ 2011 快度分布（从 3.70 到 2.14）、CMS 8 TeV $Z$ $p_T$ 分布
（从 3.65 到 1.32）以及 LHCb 8 TeV $W,Z\to \mu$
快度分布（从 1.88 到 1.37）；这些实验的非关联统计不确定度通常处于亚百分点水平。
%
随着 LHC 测量愈发精确，这一趋势很可能继续下去。


% %%%%%%%%%%%%%%%%%%%%%
% \begin{table}[t]
\begin{center}
  \scriptsize
  \renewcommand{\arraystretch}{1.10}
\begin{tabular}{|l|c|c||c|c|}
\hline
 & \multicolumn{2}{c|}{NNPDF3.1 NNLO } & \multicolumn{2}{c|}{NNPDF3.1 NLO } \\
Dataset & $\chi^2_{\rm exp}$ & $\chi^2_{t_0}$ & $\chi^2_{\rm exp}$  & $\chi^2_{t_0}$    \\
\hline
\hline
NMC    &    1.30   &   1.23   &   1.35   &   1.35  \\    
SLAC    &    0.75   &   0.73   &   1.17   &   1.12  \\    
BCDMS    &    1.21   &   1.23   &   1.17   &   1.20  \\    
CHORUS    &    1.11   &   1.14   &   1.06   &   1.09  \\    
NuTeV dimuon    &    0.82   &   0.76   &   0.87   &   0.81  \\
\hline
HERA I+II inclusive    &    1.16   &   1.26   &   1.14   &   1.29  \\    
HERA $\sigma_c^{\rm NC}$    &    1.45   &   1.46   &   1.15   &   1.15  \\    
HERA $F_2^b$   &    1.11   &   1.33   &   1.08   &   1.27  \\    
\hline
\hline
DYE866 $\sigma^d_{\rm DY}/\sigma^p_{\rm DY}$    &    0.41   &   0.41   &   0.40   &   0.40  \\    
DYE866 $\sigma^p$     &    1.43   &   1.09   &   1.05   &   0.98  \\    
DYE605    &    1.21   &   1.02   &   0.99   &   1.04  \\    
\hline
CDF $Z$ rap     &    1.48   &   1.46   &   1.62   &   1.76  \\    
CDF Run II $k_t$ jets    &    0.87   &   1.03   &   0.84   &   1.02  \\
\hline
D0 $Z$ rap    &    0.60   &   0.61   &   0.66   &   0.68  \\    
D0 $W\to e\nu$  asy    &    2.71   &   2.71   &   1.59   &   1.59  \\    
D0 $W\to \mu\nu$  asy    &    1.56   &   1.56   &   1.52   &   1.52  \\
\hline
\hline
ATLAS total   &    1.09   &   0.99   &   1.36   &   1.37  \\    
ATLAS $W,Z$ 7 TeV 2010    &    0.96   &   0.88   &   1.04   &   1.00  \\    
ATLAS HM DY 7 TeV    &    1.54   &   1.76   &   1.88   &   2.33  \\    
ATLAS low-mass DY 2011    &    0.90   &   0.99   &   0.69   &   0.78  \\    
ATLAS $W,Z$ 7 TeV 2011    &    2.14   &   2.19   &   3.70   &   4.02  \\    
ATLAS jets 2010 7 TeV     &    0.95   &   0.51   &   0.92   &   0.51  \\    
ATLAS jets 2.76 TeV     &    1.03   &   0.58   &   1.03   &   0.60  \\    
ATLAS jets 2011 7 TeV     &    1.07   &   1.51   &   1.12   &   1.57  \\    
ATLAS $Z$ $p_T$ 8 TeV $(y_{ll},M_{ll})$       &    0.93   &   1.04   &   1.17   &   1.57  \\    
ATLAS $Z$ $p_T$ 8 TeV $(p_T^{ll},M_{ll})$    &    0.94   &   1.13   &   1.76   &   2.35  \\    
ATLAS $\sigma_{tt}^{\rm tot}$     &    0.86   &   0.93   &   1.92   &   2.19  \\    
ATLAS $t\bar{t}$ rap    &    1.45   &   1.47   &   1.31   &   1.33  \\
\hline
CMS total   &    1.06   &   1.19   &   1.20   &   1.43  \\    
CMS $W$ asy 840 pb     &    0.78   &   0.78   &   0.86   &   0.86  \\    
CMS $W$ asy 4.7 fb     &    1.75   &   1.75   &   1.77   &   1.77  \\    
CMS $W+c$ tot    &    -   &   -   &   0.54   &   0.65  \\    
CMS $W+c$ ratio    &    -   &   -   &   1.91   &   1.91  \\    
CMS Drell-Yan 2D 2011    &    1.27   &   1.44   &   1.23   &   1.43  \\    
CMS $W$ rap 8 TeV    &    1.01   &   0.95   &   0.70   &   0.74  \\    
CMS jets 7 TeV 2011     &    0.84   &   0.90   &   0.84   &   0.89  \\    
CMS jets 2.76 TeV    &    1.03   &   1.31   &   1.01   &   1.28  \\    
CMS $Z$ $p_T$ 8 TeV $(p_T^{ll},M_{ll})$    &    1.32   &   1.46   &   3.65   &   5.13  \\    
CMS $\sigma_{tt}^{\rm tot}$     &    0.20   &   0.19  &   0.59   &   0.63  \\    
CMS $t\bar{t}$ rap    &    0.94   &   0.96   &   0.96   &   0.96  \\
\hline
LHCb total   &    1.47   &   1.50   &   1.61   &   1.80  \\    
LHCb $Z$ 940 pb     &    1.49   &   1.67   &   1.27   &   1.48  \\    
LHCb $Z\to ee$ 2 fb     &    1.14   &   1.16   &   1.33   &   1.44  \\    
LHCb $W,Z \to \mu$ 7 TeV    &    1.76   &   1.80   &   1.60   &   1.75  \\    
LHCb $W,Z \to \mu$ 8 TeV    &    1.37   &   1.35   &   1.88   &   2.14  \\    
\hline
\hline
Total dataset     &  {\bf  1.148}   &{\bf   1.183}   & {\bf  1.168}   & {\bf  1.246}  \\    
\hline
\end{tabular}
\caption{\small
  \label{tab:chi2tab_exp_vs_t0}
  Same as Table~\ref{tab:chi2tab_31-nlo-nnlo-30} now comparing the values of
  the  $\chi^2/N_{\rm dat}$ obtained with either  
  the experimental or the $t_0$ definition (see text), for the NNPDF3.1 NLO and NNLO fits.
}
\end{center}
\end{table}
%%%%%%%%%%%%%%%%%%%%%%%%%%%%%%%%%%%%%%%%%%%%%%%%%%%%%%%%%%%%%%%
%%%%%%%%%%%%%%%%%%%%%%%%%%%%%%%%%%%%%%%%%%%%%%%%%%%%%%%%%%%%%%%



% %%%%%%%%%%%%%%%%%%%%%

% （原文中此段比较实验 $\chi^2$ 与 $t_0$ 定义的内容被注释，保持原状不译。）

\clearpage

\subsection{部分子分布}
\label{sec:PDFcomparisons}

现在我们检视基线 NNPDF3.1 部分子分布，
并将其与 NNPDF3.0
以及 MMHT14~\cite{Harland-Lang:2014zoa}、CT14~\cite{Dulat:2015mca} 和 ABMP16~\cite{Alekhin:2017kpj} 比较。
%
NNPDF3.1 NNLO PDF 展示于图~\ref{fig:nnlopdfs}。可以看到虽然粲如今被独立参数化，它的已知精度仍高于奇异味 PDF。
在大部分实验可达的 $x$ 范围内，确定最精确的 PDF 如今是胶子，下文将更详细地讨论。


%%%%%%%%%%%%%%%%%%%%%%%%%%%%%%%%%%%%%%%%%%%%%%%%%%%%%%%%%%%%%%%%%%%%%
\begin{figure}[t]
\begin{center}
  \includegraphics[scale=0.75]{images/plots/nnpdf31nnlo-10.pdf}
  \includegraphics[scale=0.75]{images/plots/nnpdf31nnlo-1e4.pdf}
  \caption{\small NNPDF3.1 NNLO PDFs，
    分别取 $\mu^2=10~{\rm GeV}^2$（左）与 $\mu^2=10^4~{\rm GeV}^2$（右）。
    \label{fig:nnlopdfs}
  }
\end{center}
\end{figure}
%%%%%%%%%%%%%%%%%%%%%%%%%%%%%%%%%%%%%%%%%%%%%%%%%%%%%%%%%%%%%%%%%%%%%%%

 
图~\ref{fig:distances_31_vs_30} 给出 NNPDF3.1 与 NNPDF3.0 PDF 之间的距离。
按文献~\cite{Ball:2010de}的距离定义，$d\simeq 1$ 对应统计上等价的集合。比较两个 $N_{\rm rep}=100$ 副本的集合时，
距离 $d\simeq 10$ 对应以相应方差为单位的一倍标准差差异，无论对中心值还是 PDF 不确定度皆然。为清晰起见只画出总奇异味分布 $s^+=s+\bar s$ 之间的距离，而不分别显示奇异味与反奇异味。
%
我们发现对所有味道及整个 $x$ 范围，无论中心值还是 PDF 误差层面都存在重要差异。
%
距离最大的是粲——它在 NNPDF3.1 中被独立参数化而在 NNPDF3.0 中不是。除此之外，
最显著的距离出现在大 $x$ 的轻夸克分布与中等 $x$ 的奇异味上。

%%%%%%%%%%%%%%%%%%%%%%%%%%%%%%%%%%%%%%%%%%%%%%%%%%%%%%%%%%%%%%%%%%%%%
\begin{figure}[t]
\begin{center}
  \includegraphics[scale=1]{images/plots/distances_31_vs_30.pdf}
  \caption{\small NNPDF3.0 与 NNPDF3.1 NNLO PDF 集合的中心值（左）
    与不确定度（右）之间的距离，取 $Q=100$ GeV。
    注意两图 $y$ 轴标度不同。
    \label{fig:distances_31_vs_30}
  }
\end{center}
\end{figure}
%%%%%%%%%%%%%%%%%%%%%%%%%%%%%%%%%%%%%%%%%%%%%%%%%%%%%%%%%%%%%%%%%%%%%%%

图~\ref{fig:31-nnlo-vs30} 将全套 NNPDF3.1 NNLO PDF 与 NNPDF3.0 比较。
%
在 $x\lsim 0.03$ 区域 NNPDF3.1 胶子略大于 NNPDF3.0，而在更大 $x$ 处变小，
且 PDF 误差显著减小。
%
NNPDF3.1 的轻夸克与奇异味在中间 $x$ 处大于 3.0，偏差最大的是奇异味与反下夸克 PDF；
而在小 $x$ 与大 $x$ 两端两个 PDF 确定符合良好。与 NNPDF3.0 的微扰粲相比，NNPDF3.1 的最优拟合粲在中间 $x$ 区显著偏小，而在更大 $x$ 处不确定度显著增大。

对相应不确定度的详细比较见图~\ref{fig:ERR-31-nnlo-vs30}：我们比较每个 PDF 的相对不确定度，定义为其一倍标准差 PDF 不确定度与 NNPDF3.1 集合中心值之比。NNPDF3.1 不确定度要么与 NNPDF3.0 相当，要么明显更小。
唯一的重大例外是中间与大 $x$ 区的粲 PDF，其不确定度大幅增大。另一方面，
胶子 PDF 的不确定度在 NNPDF3.1 中在整个 $x$ 范围内都更小。
这是一个重要结果：由于多了一个自由参数化的 PDF，人们或许预期 NNPDF3.1 的不确定度总体更大。
唯一显著增大的不确定度恰是粲 PDF 这一事实提示：不对粲参数化可能是一种偏差来源。
中心值变化不可忽略（虽然在不确定度内相容）而自身不确定度却显著减小，
强烈提示 NNPDF3.1 比 NNPDF3.0 更准确——这与纳入拟合的大量新数据的预期一致。
粲参数化对 PDF 及其不确定度的影响将在第~\ref{sec:results-mc} 节更详细讨论，
新数据对中心值与不确定度的影响将在第~\ref{sec:disentangling} 节讨论。

%%%%%%%%%%%%%%%%%%%%%%%%%%%%%%%%%%%%%%%%%%%%%%%%%%%%%%%%%%%%%%%%%%%%%
\begin{figure}[t]
  \begin{center}
      \includegraphics[scale=0.32]{images/plots/xu-31-nnlo-vs30.pdf}
    \includegraphics[scale=0.32]{images/plots/xubar-31-nnlo-vs30.pdf}
    \includegraphics[scale=0.32]{images/plots/xd-31-nnlo-vs30.pdf}
    \includegraphics[scale=0.32]{images/plots/xdbar-31-nnlo-vs30.pdf}
  \includegraphics[scale=0.32]{images/plots/xs-31-nnlo-vs30.pdf}
  \includegraphics[scale=0.32]{images/plots/xsbar-31-nnlo-vs30.pdf}
  \includegraphics[scale=0.32]{images/plots/xc-31-nnlo-vs30.pdf}
  \includegraphics[scale=0.32]{images/plots/xg-31-nnlo-vs30.pdf}
    \caption{\small $Q=100$ GeV 处 NNPDF3.1 与 NNPDF3.0 NNLO PDF 的比较。
      自上而下依次为 up 与 antiup、down 与 antidown、strange 与 antistrange、charm 与 gluon。
    \label{fig:31-nnlo-vs30}
  }
\end{center}
\end{figure}
%%%%%%%%%%%%%%%%%%%%%%%%%%%%%%%%%%%%%%%%%%%%%%%%%%%%%%%%%%%%%%%%%%%%%%%




%%%%%%%%%%%%%%%%%%%%%%%%%%%%%%%%%%%%%%%%%%%%%%%%%%%%%%%%%%%%%%%%%%%%%
\begin{figure}[t]
  \begin{center}
    \includegraphics[scale=0.32]{images/plots/xu-ERR-31-nnlo-vs30.pdf}
    \includegraphics[scale=0.32]{images/plots/xubar-ERR-31-nnlo-vs30.pdf}
    \includegraphics[scale=0.32]{images/plots/xd-ERR-31-nnlo-vs30.pdf}
    \includegraphics[scale=0.32]{images/plots/xdbar-ERR-31-nnlo-vs30.pdf}
  \includegraphics[scale=0.32]{images/plots/xs-ERR-31-nnlo-vs30.pdf}
  \includegraphics[scale=0.32]{images/plots/xsbar-ERR-31-nnlo-vs30.pdf}
  \includegraphics[scale=0.32]{images/plots/xc-ERR-31-nnlo-vs30.pdf}
  \includegraphics[scale=0.32]{images/plots/xg-ERR-31-nnlo-vs30.pdf}
  \caption{\small $Q=100$ 处 NNPDF3.1 与 NNPDF3.0 相对
    PDF 不确定度的比较；PDF 如图~\ref{fig:31-nnlo-vs30}。
    所示不确定度均以 NNPDF3.1 中心值归一。
    \label{fig:ERR-31-nnlo-vs30}
  }
\end{center}
\end{figure}
%%%%%%%%%%%%%%%%%%%%%%%%%%%%%%%%%%%%%%%%%%%%%%%%%%%%%%%%%%%%%%%%%%%%%%%


图~\ref{fig:globalfits} 将 NNPDF3.1 PDF 与 PDF4LHC15 组合中连同 NNPDF3.0 在内的其他全局
PDF 集合即 CT14 与 MMHT14 比较。因此该比较指示了在组合中以 NNPDF3.1 替换 NNPDF3.0 的效果。
三个集合的相对不确定度在图~\ref{fig:ERR-globalfits} 中比较。把图~\ref{fig:globalfits} 与图~\ref{fig:31-nnlo-vs30} 对照可以发现一个有趣现象：
NNPDF3.1 与其他全局拟合之间差异模式的多个方面与 NNPDF3.1 和 NNPDF3.0 之间的情形相似，
因而很可能有类似的起源。这对粲与胶子尤为明显。
在与希格斯产生相关的 $x\lsim 0.03$ 区域，
三者的胶子仍然相当一致；但现在 NNPDF3.1 位于一倍标准差范围的上沿，即该区域的 NNPDF3.1 胶子偏高。
%
在大 $x$ 处，NNPDF3.1 胶子相对 MMHT14 与 CT14 则被压低。正如我们将在第~\ref{sec:results-mc}、\ref{sec:disentangling} 节所示，
该增强是粲参数化的结果；而如第~\ref{sec:impacttop} 节所示，大 $x$ 的压低则是纳入 8~TeV 顶夸克微分数据的直接后果。
如今 NNPDF3.1 胶子 PDF 的不确定度已明显小于 CT14 或 MMHT14 任一家。

对夸克 PDF，up 与 down 在整个 $x$ 范围内符合良好。antidown PDF 的符合则处于边缘状态：
$x\lsim 0.1$ 时 NNPDF3.1 高于 MMHT14 与 CT14，更大 $x$ 时低于二者。质子的奇异味占比在 NNPDF3.1 中大于 CT14 与
MMHT14，且 PDF 不确定度要小得多。
与 CT14、MMHT14 的微扰生成粲相比，NNPDF3.1 的最优拟合粲在中间 $x$ 受到压低，但在大 $x$ 具有大得多的不确定度——这自然应如此，差异显然可追溯到 NNPDF3.1 中粲被自由参数化这一事实。

%%%%%%%%%%%%%%%%%%%%%%%%%%%%%%%%%%%%%%%%%%%%%%%%%%%%%%%%%%%%%%%%%%%%%
\begin{figure}[t]
  \begin{center}
 \includegraphics[scale=0.32]{images/plots/xu-31-nnlo-globalfits.pdf}
    \includegraphics[scale=0.32]{images/plots/xubar-31-nnlo-globalfits.pdf}
    \includegraphics[scale=0.32]{images/plots/xd-31-nnlo-globalfits.pdf}
    \includegraphics[scale=0.32]{images/plots/xdbar-31-nnlo-globalfits.pdf}
  \includegraphics[scale=0.32]{images/plots/xs-31-nnlo-globalfits.pdf}
  \includegraphics[scale=0.32]{images/plots/xsbar-31-nnlo-globalfits.pdf}
  \includegraphics[scale=0.32]{images/plots/xc-31-nnlo-globalfits.pdf}
  \includegraphics[scale=0.32]{images/plots/xg-31-nnlo-globalfits.pdf}
  \caption{\small NNPDF3.1、CT14 与 MMHT2014 NNLO PDF 的比较。
    比较在 $Q=100$ GeV 进行，结果以 NNPDF3.1 中心值归一；PDF 如图~\ref{fig:31-nnlo-vs30}。
    \label{fig:globalfits}
  }
\end{center}
\end{figure}
%%%%%%%%%%%%%%%%%%%%%%%%%%%%%%%%%%%%%%%%%%%%%%%%%%%%%%%%%%%%%%%%%%%%%%%

%%%%%%%%%%%%%%%%%%%%%%%%%%%%%%%%%%%%%%%%%%%%%%%%%%%%%%%%%%%%%%%%%%%%%
\begin{figure}[t]
  \begin{center}
 \includegraphics[scale=0.32]{images/plots/xu-ERR-31-nnlo-globalfits.pdf}
    \includegraphics[scale=0.32]{images/plots/xubar-ERR-31-nnlo-globalfits.pdf}
    \includegraphics[scale=0.32]{images/plots/xd-ERR-31-nnlo-globalfits.pdf}
    \includegraphics[scale=0.32]{images/plots/xdbar-ERR-31-nnlo-globalfits.pdf}
  \includegraphics[scale=0.32]{images/plots/xs-ERR-31-nnlo-globalfits.pdf}
  \includegraphics[scale=0.32]{images/plots/xsbar-ERR-31-nnlo-globalfits.pdf}
  \includegraphics[scale=0.32]{images/plots/xc-ERR-31-nnlo-globalfits.pdf}
  \includegraphics[scale=0.32]{images/plots/xg-ERR-31-nnlo-globalfits.pdf}
  \caption{\small $Q=100$ 处 NNPDF3.1、CT14 与 MMHT2014 相对
    PDF 不确定度的比较；PDF 如图~\ref{fig:globalfits}。
    \label{fig:ERR-globalfits}
  }
\end{center}
\end{figure}
%%%%%%%%%%%%%%%%%%%%%%%%%%%%%%%%%%%%%%%%%%%%%%%%%%%%%%%%%%%%%%%%%%%%%%%



%%%%%%%%%%%%%%%%%%%%%%%%%%%%%%%%%%%%%%%%%%%%%%%%%%%%%%%%%%%%%%%%%%%%%
\begin{figure}[t]
  \begin{center}
    \includegraphics[scale=0.32]{images/plots/xu-abmp16.pdf}
    \includegraphics[scale=0.32]{images/plots/xubar-abmp16.pdf}
    \includegraphics[scale=0.32]{images/plots/xd-abmp16.pdf}
    \includegraphics[scale=0.32]{images/plots/xdbar-abmp16.pdf}
  \includegraphics[scale=0.32]{images/plots/xs-abmp16.pdf}
  \includegraphics[scale=0.32]{images/plots/xsbar-abmp16.pdf}
  \includegraphics[scale=0.32]{images/plots/xc-abmp16.pdf}
  \includegraphics[scale=0.32]{images/plots/xg-abmp16.pdf}
  \caption{\small 同图~\ref{fig:globalfits}，但此处与 ABMP16 NNLO $n_f=5$
    集合比较，分别用其默认 $\alpha_s(m_Z)=0.1147$ 以及 $\alpha_s(m_Z)=0.118$。
    \label{fig:abmp16}
  }
\end{center}
\end{figure}
%%%%%%%%%%%%%%%%%%%%%%%%%%%%%%%%%%%%%%%%%%%%%%%%%%%%%%%%%%%%%%%%%%%%%%%

最后，图~\ref{fig:abmp16} 将 NNPDF3.1
与近期的 ABMP16 集合比较。该集合在各种固定味数方案下发布。由于我们在 $Q^2=10^4$ GeV$^2$ 标度做比较，
选用 $n_f=5$ NNLO ABMP16 集合，既用其默认值 $\alpha_s(m_Z)=0.1147$ 也用 $\alpha_s(m_Z)=0.118$。
采用共同值 $\alpha_s(m_Z)=0.118$ 时，胶子 PDF 总体符合合理，
除了大 $x$ 处 ABMP16 低于
NNPDF3.1。
%
轻夸克的差别较大：
ABMP16 的 up 分布在大 $x$ 高于 NNPDF3.1，而 down 夸克在整个 $x$ 范围偏低。
奇异味 PDF 差别最大；不过对照图~\ref{fig:globalfits} 可以清楚看到 ABMP16 相对 MMHT14 与 CT14 也有类似大小的偏离。
总体而言 ABMP16
集合的不确定度比 NNPDF3.1 小得多。
这在奇异味上尤为醒目，两者不确定度差异格外突出。与之形成对比的是 CT14 与 MMHT14，
它们在整个数据区的不确定度与 NNPDF3.1 定性相似。
ABMP16 不确定度如此之小的原因可追溯到其过于受限的参数化，以及该集合虽用 Hessian 方法产生，
却不像 MMHT14 与 CT14 那样设置容忍度（tolerance）（见文献~\cite{Harland-Lang:2014zoa,Dulat:2015mca}）。



\clearpage

\clearpage

% 对粲质量的稳定性
\subsection{方法学改进：粲的参数化}
\label{sec:results-mc}

NNPDF3.1 相对 NNPDF3.0 最主要的方法学改进在于：粲 PDF 如今与轻夸克及奇异味夸克 PDF 以同样方式参数化。
为量化这一变化的影响，我们重复了 NNPDF3.1 分析，但对粲的处理沿用以往所有 NNPDF 确定，
即完全通过 NLO 或 NNLO 匹配条件微扰生成。
%%%%%%%%%%%%%%%%%%%%
\begin{table}[h]
\centering
\includegraphics[page=1,width=0.88\textwidth]{images/D.pdf}
\caption{\small 同表~\ref{tab:chi2tab_31-nlo-nnlo-30}，但此处比较默认 NNPDF3.1 NNLO/NLO 集合与微扰生成粲的变体。对 HERA $\sigma_c^{\rm NC}$，括号中的数值对应施加了表~\ref{eq:tablekincuts} 中 NNLO 拟合粲（FC）切割的数据子集。
\label{tab:chi2tab_31-nlo-nnlo-30-pc}}
\end{table}


%%%%%%%%%%%%%%%%%%%%%

表~\ref{tab:chi2tab_31-nlo-nnlo-30-pc} 给出粲在 NLO 与 NNLO 微扰生成时的 $\chi^2/N_{\rm dat}$ 数值。毫不意外，当粲不被独立参数化时拟合质量变差。
这是朴素预期的结果，
因为微扰粲对拟合施加了约束，从而减少了自由参数的数目。

 然而，有趣的是：采用微扰粲时，单举 HERA 数据（1306 个数据点）的拟合质量
 从 NLO 到 NNLO 显著恶化，而当粲独立参数化时保持稳定。关于粲结构函数数据，请注意如上文第~\ref{sec:datadis} 节所讨论，
 当粲被独立参数化时，NNLO 下会对 HERA $\sigma_c^{\rm NC}$ 数据施加额外切割。为使比较自洽，
 表~\ref{tab:chi2tab_31-nlo-nnlo-30-pc} 中括号内给出的是对该切割后幸存的 47 点中的 37 点
 在其他各种情形下计算的 $\chi^2/N_{\rm dat}$。因此对该数据而言，
 微扰粲与参数化粲的拟合质量相近，NLO 与 NNLO 也相近（NNLO 略差，幅度与统计涨落相容）。
 参数化粲之后从 NLO 到 NNLO 拟合质量不再恶化这一事实提示：
 这化解了 NNLO 下微扰粲方案中 HERA 数据与强子对撞机数据之间的张力。同样，纯微扰粲导致 BCDMS、NMC 特别是 NuTeV 双 $\mu$ 截面在 NNLO 明显变差。
 这可追溯到：独立参数化粲对于调和 HERA 数据与 ATLAS $W,Z$ 2011 快度分布对质子奇异味含量所加约束至关重要。

图~\ref{fig:31-nnlo-fitted-vs-pch} 直接比较了参数化粲与微扰粲两种 PDF。
当粲被独立参数化时，轻夸克 PDF 与胶子在 $x\gsim0.003$ 一般被增强、更小 $x$ 处被压低。
差异最大的是 up 夸克，而奇异味与胶子分布更为稳定。最优拟合粲分布与其微扰生成对应物形状截然不同且不确定度显著更大。
如文献~\cite{Ball:2016neh}所论，该形状很可能与高微扰阶生成的粲 PDF 相容。

%%%%%%%%%%%%%%%%%%%%%%%%%%%%%%%%%%%%%%%%%%%%%%%%%%%%%%%%%%%%%%%%%%%%%
\begin{figure}[t]
  \begin{center}
    \includegraphics[scale=0.32]{images/plots/xu-31-nnlo-fitted-vs-pch.pdf}
    \includegraphics[scale=0.32]{images/plots/xd-31-nnlo-fitted-vs-pch.pdf}
    \includegraphics[scale=0.32]{images/plots/xubar-31-nnlo-fitted-vs-pch.pdf}
    \includegraphics[scale=0.32]{images/plots/xdbar-31-nnlo-fitted-vs-pch.pdf}
    \includegraphics[scale=0.32]{images/plots/xs-31-nnlo-fitted-vs-pch.pdf}
    \includegraphics[scale=0.32]{images/plots/xsbar-31-nnlo-fitted-vs-pch.pdf}
    \includegraphics[scale=0.32]{images/plots/xc-31-nnlo-fitted-vs-pch.pdf}
    \includegraphics[scale=0.32]{images/plots/xg-31-nnlo-fitted-vs-pch.pdf}
  \caption{\small 
    NNPDF3.1 NNLO PDF 与粲完全微扰生成（其余一切不变）变体的比较。
    \label{fig:31-nnlo-fitted-vs-pch}
  }
\end{center}
\end{figure}
%%%%%%%%%%%%%%%%%%%%%%%%%%%%%%%%%%%%%%%%%%%%%%%%%%%%%%%%%%%%%%%%%%%%%%%

图~\ref{fig:31-nnlo-fitted-vs-pch-err} 直接比较 PDF 不确定度。值得注意的是，
当粲被独立参数化时除粲以外的不确定度基本不变，海夸克 PDF 不确定度仅在 $10^{-3}\lsim x \lsim 10^{-2}$ 有轻微增大。
胶子的不确定度几乎完全不受影响。粲被独立参数化时其 PDF 不确定度与其他海夸克 PDF 相当；
而微扰生成粲的不确定度跟随胶子，因而小得多。

%%%%%%%%%%%%%%%%%%%%%%%%%%%%%%%%%%%%%%%%%%%%%%%%%%%%%%%%%%%%%%%%%%%%%
\begin{figure}[t]
  \begin{center}
    \includegraphics[scale=0.32]{images/plots/xu-ERR-31-nnlo-fitted-vs-pch.pdf}
    \includegraphics[scale=0.32]{images/plots/xd-ERR-31-nnlo-fitted-vs-pch.pdf}
    \includegraphics[scale=0.32]{images/plots/xubar-ERR-31-nnlo-fitted-vs-pch.pdf}
    \includegraphics[scale=0.32]{images/plots/xdbar-ERR-31-nnlo-fitted-vs-pch.pdf}
    \includegraphics[scale=0.32]{images/plots/xs-ERR-31-nnlo-fitted-vs-pch.pdf}
    \includegraphics[scale=0.32]{images/plots/xsbar-ERR-31-nnlo-fitted-vs-pch.pdf}
    \includegraphics[scale=0.32]{images/plots/xc-ERR-31-nnlo-fitted-vs-pch.pdf}
    \includegraphics[scale=0.32]{images/plots/xg-ERR-31-nnlo-fitted-vs-pch.pdf}
  \caption{\small 
    NNPDF3.1 NNLO 的一倍标准差分数 PDF 不确定度与粲微扰生成版本
    （其余一切不变）的比较。
    相应的 PDF 比较图示于
    图~\ref{fig:31-nnlo-fitted-vs-pch}。
    \label{fig:31-nnlo-fitted-vs-pch-err}
  }
\end{center}
\end{figure}
%%%%%%%%%%%%%%%%%%%%%%%%%%%%%%%%%%%%%%%%%%%%%%%%%%%%%%%%%%%%%%%%%%%%%%%

此前文献~\cite{Ball:2016neh}曾对参数化粲与微扰粲确定的 PDF 进行过比较，结论是：
将粲参数化并由数据确定会大大降低对粲质量的依赖，从而在计入粲质量不确定度时降低整体 PDF 不确定度。
如前所述，NNPDF3.1 PDF 采用希格斯截面工作组~\cite{deFlorian:2016spz}推荐的重夸克极点质量及其不确定度。
对粲而言即 $m_c^{\rm pole}=1.51 \pm 0.13\,{\rm GeV}$。
%
为估计该不确定度的影响，我们产生了 $m_c^{\rm pole}=1.38$~GeV 与 $m_c^{\rm pole}=1.64$~GeV 的
NNPDF3.1 NNLO 集合。
结果对若干代表性 PDF 示于图~\ref{fig:nnlo-fc-mcdep-1}，同时给出默认 NNPDF3.1 与微扰粲版本。
显然，当粲被独立参数化后，微扰生成粲时发现的粲 PDF 对 $m_c$ 的极强依赖几乎完全消失。
胶子始终相当稳定，而微扰生成粲对 $m_c$ 的依赖会传播到轻夸克分布；参数化粲使其得到显著稳定。
事实上若粲为微扰生成，粲质量变动一倍标准差引起的 up、down 夸克分布位移可与（虽略小于）
PDF 不确定度相当。
当粲独立参数化时该依赖大大减弱。
采用参数化粲后，高标度对撞机观测量变得本质上与粲质量无关，与解耦论证的预期一致。










%%%%%%%%%%%%%%%%%%%%%%%%%%%%%%%%%%%%%%%%%%%%%%%%%%%%%%%%%%%%%%%%%%%%%
\begin{figure}[t]
\begin{center}
  \includegraphics[scale=0.36]{images/plots/xc-31-nnlo-fc-mcdep.pdf}
  \includegraphics[scale=0.36]{images/plots/xc-31-nnlo-pc-mcdep.pdf}
  \includegraphics[scale=0.36]{images/plots/xg-31-nnlo-fc-mcdep.pdf}
  \includegraphics[scale=0.36]{images/plots/xg-31-nnlo-pc-mcdep.pdf}
  \includegraphics[scale=0.36]{images/plots/xu-31-nnlo-fc-mcdep.pdf}
  \includegraphics[scale=0.36]{images/plots/xu-31-nnlo-pc-mcdep.pdf}
  \includegraphics[scale=0.36]{images/plots/xd-31-nnlo-fc-mcdep.pdf}
  \includegraphics[scale=0.36]{images/plots/xd-31-nnlo-pc-mcdep.pdf}
  \caption{\small NNPDF3.1 NNLO PDF 对粲质量的依赖。参数化粲（左）
    与微扰粲（右）均予给出，（自上而下）依次为 charm、gluon、up、down。
        \label{fig:nnlo-fc-mcdep-1}
  }
\end{center}
\end{figure}
%%%%%%%%%%%%%%%%%%%%%%%%%%%%%%%%%%%%%%%%%%%%%%%%%%%%%%%%%%%%%%%%%%%%%%%









\clearpage

% 理论不确定度的讨论
\subsection{理论不确定度}
\label{sec:thunc}

进入 PDF4LHC15 组合的全局 PDF 集合，其 PDF 不确定度
只包含由实验数据传播的不确定度与方法学带来的不确定度。
后者可通过封闭性检验加以控制。然而还存在源自 PDF 确定所用理论的进一步不确定度来源，这里作简要评估。它们可分为两大类：
\begin{itemize}
\item 缺失高阶不确定度（MHOU），源于 PDF 确定所用理论在某一固定阶（LO、NLO 或 NNLO）处截断 QCD 微扰展开。
\item 参数不确定度，来自 PDF 确定所用理论参数取值的不确定度：主要是 $\alpha_s(m_Z)$ 与 $m_c^{\rm pole}$ 的取值。
\end{itemize}

对 MHOU 的完整评估是一个悬而未决的问题，留待未来研究。目前可通过研究结果的微扰稳定性获得初步评估。
图~\ref{fig:distances_perturbativeorder} 给出 $Q=100$ GeV 处
 LO 与 NLO 集合之间、以及 NLO 与 NNLO 集合之间所有 PDF 的距离。随后在
 图~\ref{fig:pdfcomp-pertorder} 中直接比较若干 LO、NLO 与 NNLO PDF。
LO 与 NLO 集合的差异非常大，无论中心值还是不确定度皆是如此；后者在 LO 很大是因为
拟合质量很差。
夸克 PDF 的位移可高达两倍标准差（$d\simeq 20$），而小 $x$ 胶子在 LO 与 NLO 间截然不同，原因在于夸克到胶子劈裂函数的奇异小 $x$ 行为只在 NLO 才开始出现，并且胶子 initiated 的 DIS 与 DY 过程在 LO 为零。
另一方面，从 NLO 到 NNLO，
PDF 不确定度基本不受影响。中心值也相当稳定：最大的位移出现在大 $x$ 胶子、大 $x$ down 夸克与小 $x$ 胶子，仍不超过一倍标准差水平。

%%%%%%%%%%%%%%%%%%%%%%%%%%%%%%%%%%%%%%%%%%%%%%%%%%%%%%%%%%%%%%%%%%%%%
\begin{figure}[t]
\begin{center}
  \includegraphics[scale=1]{images/plots/distances_lo_vs_nlo.pdf}
  \includegraphics[scale=1]{images/plots/distances_nlo_vs_nnlo.pdf}
  \caption{\small $Q=100$ GeV 处 LO 与 NLO（上）、NLO 与 NNLO（下）的 NNPDF3.1
    PDF 之间的距离。注意两图 $y$ 轴标度不同。
    \label{fig:distances_perturbativeorder}
  }
\end{center}
\end{figure}
%%%%%%%%%%%%%%%%%%%%%%%%%%%%%%%%%%%%%%%%%%%%%%%%%%%%%%%%%%%%%%%%%%%%%%%


%%%%%%%%%%%%%%%%%%%%%%%%%%%%%%%%%%%%%%%%%%%%%%%%%%%%%%%%%%%%%%%%%%%%%
\begin{figure}[t]
\begin{center}
  \includegraphics[scale=0.33]{images/plots/xg-pertorder.pdf}
  \includegraphics[scale=0.33]{images/plots/xu-pertorder.pdf}
  \includegraphics[scale=0.33]{images/plots/xdbar-pertorder.pdf}
  \includegraphics[scale=0.33]{images/plots/xsp-pertorder.pdf}
  \caption{\small 若干 LO、NLO 与 NNPDF3.1 NNLO PDF 的比较：gluon 与 up（上）、
    antidown 与 total strangeness（下）。所有结果均在 $Q=100$ GeV 给出并以 NNLO 中心值归一。}
    \label{fig:pdfcomp-pertorder}
\end{center}
\end{figure}
%%%%%%%%%%%%%%%%%%%%%%%%%%%%%%%%%%%%%%%%%%%%%%%%%%%%%%%%%%%%%%%%%%%%%%%

通过计算 NLO 与 NNLO NNPDF3.1 PDF 中心值之间的位移可获得 MHOU 的定量估计。
结果对若干 PDF 组合示于图~\ref{fig:THerror-shifts}。图中位移已对称化，
并与 NLO 标准 PDF 不确定度比较。对夸克单态 $\Sigma$，在 $x\lesssim 10^{-3}$ 位移大于 PDF 不确定度，
而对单独味道（如图所示的两个夸克分布）则更小。
这提示在当前精度水平下，对单独夸克味道与 NNLO 胶子而言 MHOU 可以合理地忽略。
然而对特定组合（如小 $x$ 单态），鉴于 NLO 时 MHOU 已超过 PDF 不确定度，
即使在 NNLO 能否忽略 MHOU 也不清楚。

%
%%%%%%%%%%%%%%%%%%%%%%%%%%%%%%%%%%%%%%%%%%%%%%%%%%%%%%%%%%%%%%%%%%%%%
\begin{figure}[t]
\begin{center}
  \includegraphics[scale=0.36]{images/plots/xsinglet-THerror.pdf}
  \includegraphics[scale=0.36]{images/plots/xg-THerror.pdf}
  \includegraphics[scale=0.36]{images/plots/xu-THerror.pdf}
  \includegraphics[scale=0.36]{images/plots/xdbar-THerror.pdf}
  \caption{\small NLO PDF 不确定度与 NLO、NNLO PDF 位移的比较。
    所有结果都以对 NLO PDF 之比给出，取 $Q=100$~GeV。位移已对称化。
    上排为 singlet、gluon；下排为 up 与 antidown。
    \label{fig:THerror-shifts}
  }
\end{center}
\end{figure}
%%%%%%%%%%%%%%%%%%%%%%%%%%%%%%%%%%%%%%%%%%%%%%%%%%%%%%%%%%%%%%%%%%%%%%%

最后转向参数不确定度。如第~\ref{sec:results-mc} 节所讨论，
参数化粲几乎完全消除了 PDF 对粲质量的依赖。
对底夸克质量的依赖很小，底 PDF 本身除外~\cite{Ball:2011mu,Harland-Lang:2015qea}。
因此唯一重要的残余参数不确定度来自强耦合取值。
该不确定度通常与 PDF 不确定度一并计入；为自洽起见需要用不同 $\alpha_s$ 中心值产生的 PDF 集合（例如见
文献~\cite{Butterworth:2015oua}）。我们已在 $0.108\le\alpha_s(m_Z)\le0.124$ 范围内确定了
NNPDF3.1 NLO 与 NNLO PDF（见第~\ref{sec:delivery} 节）。

%%%%%%%%%%%%%%%%%%%%%%%%%%%%%%%%%%%%%%%%
\begin{figure}[t]
  \begin{center}
  \includegraphics[scale=0.36]{images/plots/xg-31-nlo-fc-asdep.pdf}
  \includegraphics[scale=0.36]{images/plots/xu-31-nlo-fc-asdep.pdf}
  \includegraphics[scale=0.36]{images/plots/xg-31-nnlo-fc-asdep.pdf}
  \includegraphics[scale=0.36]{images/plots/xu-31-nnlo-fc-asdep.pdf}
  \caption{\small 
    \label{fig:PDFasvariations} NNPDF3.1 NLO（上）与 NNLO（下）PDF 对
    $\alpha_s$ 取值的依赖。gluon（左）与 up quark（right) 在 $Q=100$~GeV 给出并以中心值归一。
  }
\end{center}
\end{figure}
%%%%%%%%%%%%%%%%%%%%%%%%%%%%%%%%%%%%%%%%

图~\ref{fig:PDFasvariations} 比较了 $\alpha_s(m_Z)$ 在中心值上下变动 $\Delta\alpha_s=\pm0.002$ 时的 up 与 gluon PDF。
众所周知，
在小 $x$ 与中等 $x$ 处胶子与 $\alpha_s(m_Z)$ 反关联，而在大 $x$ 正关联。
夸克 PDF 对 $\alpha_s$ 的依赖温和得多：小 $x$ 处正关联，大 $x$ 处几乎没有依赖。




\clearpage
